%% 297_HLV_Untermenge_Kandidat_En_ch.tex
%% FFGFT -- Dok. 297: HLV as possible subset -- candidate note
%% Author: Johann Pascher
%% Date: July 2026
%% Language: English

\section*{Doc.~297 \quad HLV as a possible subset of FFGFT ---
Candidate note}

\subsection*{Starting point}

In the course of the bridge analysis between FFGFT and HLV, a finding
emerged that is recorded here as a directed candidate note. It is not a
closed derivation; it names the next clean step (P35).

\subsection*{Three steps of retrospect}

\textbf{Step 1 --- original blockade.} The $C_3$-in-$A_5$ geometry of the
HLV framework (icosahedral projector, $\tau = \phi$, $\theta = 2/9$ as
output) was not directly visible on $T^4/\mathbb{Z}_3$. A direct geometric
identification was blocked: $T^4$ is compact 4-dimensional, $A_5$ is the
symmetry group of the icosahedron in 3D --- the dimensions do not match
straightforwardly.

\textbf{Step 2 --- Hilbert-space transformation opens the path.} The
lossless bijection $H = L^2(T^4)\otimes\mathbb{C}^3$ (Doc.~230/282)
translates $T^4/\mathbb{Z}_3$ into an operator algebra on the
$\mathbb{C}^3$ fibre. In $\mathbb{C}^3$, $C_3 < A_5$ is naturally
accessible: the $\mathbb{Z}_3$ circulant diagonalisation on the fibre
outputs $\theta = 2/9$ as a phase (Doc.~291), and the $C_3$-in-$A_5$
geometry is demonstrable as a parameter-free invariant on the same fibre
(Doc.~293). The transformation made visible what was geometrically hidden
directly.

\textbf{Step 3 --- consequence.} If FFGFT contains $C_3 < A_5$ via the
Hilbert-space transformation, and HLV is built on $C_3 < A_5$ with
$\tau = \phi$, then HLV's geometric core could be a \emph{subset} of the
FFGFT structure --- not a peer framework, but a specific projection of
FFGFT's $\mathbb{C}^3$-fibre geometry.

\subsection*{What is established (from available documents)}

\begin{itemize}
  \item FFGFT contains $C_3 < A_5$ and $\theta = 2/9$ via the
    Hilbert-space transformation --- documented in Doc.~291 (circulant
    phase) and Doc.~293 (parameter-free $C_3$-in-$A_5$ invariant, second
    witness).
  \item HLV's geometric object (the $\phi$-icosahedral projector) carries
    the same $C_3 < A_5$ --- one-sidedly verified (FFGFT $\to$ HLV,
    Doc.~294).
  \item The verification direction is asymmetric: FFGFT derives $\theta =
    2/9$; HLV books it as a candidate overlap without its own derivation
    --- confirmed by inspection of his available documents (Stage~19
    manuscript, HLV-Consciousness-Bridge1, Runs~H--J): $\theta = 2/9$
    appears in none of HLV's internal derivations.
  \item His $H_2$ operator (Run~H, \texttt{skew\_J}) contains $1/\phi$
    and $1/\phi^2$, but no $C_3$-in-$A_5$ output and no $\theta = 2/9$
    as a result. Run~H is moreover \emph{negative/underdetermined} (does
    not survive the $u$-shuffle specificity test).
\end{itemize}

\subsection*{What is not shown (open obligation)}

That HLV's $\phi$-icosahedral layer is \emph{fully} contained in the
FFGFT $\mathbb{C}^3$ fibre requires an explicit \emph{embedding map}
\begin{equation}
  \iota\colon \text{HLV ground object} \hookrightarrow
  L^2(T^4)\otimes\mathbb{C}^3,
\end{equation}
which has not yet been constructed. \enquote{Subset} is the structurally
plausible conjecture --- not a proven theorem.

\subsection*{Next clean step}

Construct or falsify the embedding map $\iota$:

\begin{enumerate}
  \item Show that HLV's $\phi$-icosahedral projector is representable as
    an operator on FFGFT's $\mathbb{C}^3$ fibre --- i.e.\ as an element
    of $\mathrm{End}(\mathbb{C}^3)$ with $C_3$ invariance.
  \item Check whether HLV's spectral data (Run~J, native-6D candidate)
    are compatible with the discrete spectrum of the $T^4$
    connection-Laplacian (Doc.~283) --- or whether they deviate and thus
    block the embedding.
  \item If $\iota$ exists: HLV is a specific projection of FFGFT --- its
    geometric layer is a subset, while its dynamical objects (Spiral-Time,
    $U_3$) remain independent and unembedded.
  \item If $\iota$ cannot be constructed: both frameworks share a common
    object ($C_3 < A_5$, $\theta = 2/9$) but are structurally independent
    --- no subset relation.
\end{enumerate}

\subsection*{Concrete embedding candidate: the Markov limit and below}

The question of whether Marcel's Archimedean spiral time finds a
structurally defined place in FFGFT can be answered precisely --- not as
an analogy but as a limit. And the analysis reveals that the object is
even weaker than initially assumed.

\textbf{Archimedean = Markov limit --- but U3 is weaker still.} The FFGFT
memory kernel (Fourier transform of the discrete spectral density $\sum_k
g_k^2\,\delta(\omega-\omega_k)$, Doc.~283) becomes Archimedean exactly
when the couplings $g_k \to 0$ --- the \emph{Markov limit} of open
quantum systems. In this limit the kernel still accumulates without bound
on $\mathbb{R}$. Marcel's $U_3$ (Eq.~7 of his Consciousness document)
goes further: the window $M_m$ is finite and fixed, older values actively
drop out (\emph{sliding window}). This yields a \textbf{hard-reset kernel}
--- a memory kernel with finite support:
\begin{equation}
  K(t-\tau) = 0 \quad \text{for } t-\tau > M_m\cdot\Delta t.
\end{equation}
This is not merely Archimedean --- it is \emph{periodically bounded}.
$U_3$ does not accumulate without bound; it closes after $M_m$ steps.

\textbf{Does the spiral close before the first turn?} In FFGFT's frozen
configuration (Doc.~295), the winding closes after exactly 75 turns
($1/(100\xi_0) = 75$) --- the smallest coherent unit of the recursion.
If Marcel's window width $M_m < 75$ steps (in the winding coordinate),
then $U_3$ does not even accumulate one complete turn. There is no closed
turn, no geometric spiral --- only a sliding memory window that renews
with each step. The object lies \textbf{below the smallest coherent unit}
of the FFGFT recursion.

\textbf{Refined embedding candidate.} The concrete location of $U_3$ in
FFGFT is therefore not the Markov limit (which still accumulates on
$\mathbb{R}$) but the \emph{hard-reset kernel below one turn}:
\begin{equation}
  \iota\colon U_3 \;\hookrightarrow\;
  \bigl\{K \in L^2(T^4)\otimes\mathbb{C}^3 \;\big|\;
  K(t-\tau) = 0 \text{ for } t-\tau > M_m\cdot\Delta t,\;
  M_m < 75\bigr\}.
\end{equation}
This is present in FFGFT --- but as the \textbf{degenerate boundary case},
not as a regular configuration. The hierarchy is therefore clear:
\begin{equation}
  U_3 \;\subset\; \text{Markov limit} \;\subset\;
  \text{Archimedean} \;\subset\; \text{log spiral (}e\text{)}
  \;=\; \text{FFGFT kernel (Doc.~295/283)}.
\end{equation}
Each level is more information-poor than the next. $U_3$ is the poorest.

\subsection*{Falsification condition}

The embedding $\iota$ is falsifiable: if Marcel's $U_3$ carries
\emph{more} in his data than the hard-reset kernel --- i.e.\ if non-Markov
signatures (revivals, BLP backflow, oscillatory kernel, or correlations
beyond $M_m$ steps) are detectable --- then $U_3$ exceeds the boundary
case and a deeper embedding is needed.

Marcel's own discipline (U3 ablation test, Table~4 of his Consciousness
document) is exactly the right test: if $U_3$ does not exceed the Markov
baseline, the embedding in the boundary case is correct; if it does, a
richer layer is needed. The programme is clean --- he only needs to run it
to the end.

\subsection*{Two HLV spiral times --- two different embedding candidates}

Doc.~295 §9 records a distinction that is also load-bearing for the
embedding question of Doc.~297: \enquote{Spiral time} in HLV denotes not
one object but two.

\textbf{First object --- the geometric time} from the factorisation
$T^7 = T^4 \times T^3$ (Doc.~285): an $A_5$ singlet on the $S^1$ factor,
decoupled and separable. This limit corresponds in the
K\textsubscript{frak} recursion to the fixed point $\xi \to 0$ --- the
defect vanishes, the winding closes after 75 turns, the pitch becomes
constant. This is the \emph{geometric} boundary of the FFGFT spiral.

\textbf{Second object --- the temporal-dynamic spiral time}, the object
HLV \emph{calls} by that name: the triadic operator
$\Psi(t) = (t, \phi(t), \chi(t))$ with a memory channel, without $A_5$
and without an $S^1$ factor. This operator is introduced in Krueger et
al.\ (2026, \texttt{doi:10.47852/bonviewMEDIN62029043}) as a deterministic
phase-memory decomposition and implemented as a numerical operator in the
HLV G-Lattice Runs~H--J (Zenodo, DOI not yet published). In the
Consciousness document (Krueger, 2026, v0.2, not yet on Zenodo) it appears
as $\Psi(t) = (U_1, U_2, U_3)$ with the explicit implementation Eq.~7 ---
a bounded window accumulator.

\textbf{Two different embedding candidates in FFGFT:}

The \emph{geometric} object ($A_5$ singlet, $S^1$, decoupled) finds its
natural place at the $\xi \to 0$ fixed point of the recursion --- the
geometrically separable, information-poorest configuration of compact time.
It is geometrically richer than $U_3$, but carries no ratio $e$ (Doc.~295
§10: blocked on the $e$-criterion).

The \emph{temporal-dynamic} object ($U_3$, hard-reset kernel, $M_m < 75$)
finds its natural place in the Markov limit of the FFGFT memory kernel
(Doc.~283, $g_k \to 0$) --- below the smallest coherent unit of the
recursion. It is the more information-poor of the two.

The hierarchy applies to the temporal-dynamic object:
\[
  U_3 \;\subset\; \text{Markov limit} \;\subset\;
  \text{Archimedean} \;\subset\; \text{log spiral}~(e)
  = \text{FFGFT kernel}.
\]
The geometric singlet lies to the side of this hierarchy --- it is
Archimedean ($\xi \to 0$, constant pitch), but from geometric
factorisation, not from dynamic accumulation. Both boundary cases meet in
the Archimedean regime, but they are different paths to it.

Regardless of whether $U_3$ is represented linearly, in polar coordinates,
or as a spiral, the structure of the dependence remains the same: finite
short-term memory with a fixed horizon $M_m$, without accumulation. The
spiral representation is a \emph{choice of coordinates}, not a structural
property of the object --- in every representation $U_3$ remains a
hard-reset kernel with finite support $M_m < 75$. The geometric richness
lies in the representation, not in the object. This is the same line as
P35 (\enquote{translation is not justification}) applied in the other
direction: a richer representation does not make the object richer.

The deeper difference lies elsewhere: \textbf{accumulation is the actual
memory achievement}. $U_3$ does not accumulate --- it replaces. What was
before $M_m$ steps is irretrievably gone; no past state acts on the future
beyond the window. That is a short-term buffer, not memory in the
structural sense.

The K\textsubscript{frak} recursion is the opposite:
$\xi_{n+1} = \xi_n(1 - 100\,\xi_n)$ accumulates each step into the next
--- not as an explicit store, but as \emph{structural imprinting}: the
present state $\xi_n$ carries the entire history of all preceding steps
within itself. No step is forgotten, no window erases. The past is encoded
in the object itself, not in a separate memory. This is memory in the
deepest sense --- not as storage but as a generative structure that carries
its own history and from it produces the log spiral with ratio $e$
(Doc.~295).

That is precisely what Marcel's consciousness framework requires (Section~1
of his Consciousness document: \enquote{memory-bearing, recursively
integrated}) --- but what $U_3$, by construction, does not deliver. The
recursion would be the right carrier. $U_3$ is the degenerate boundary
case below it.

\subsection*{The window as reduction --- in the sense of Doc.~296}

The window $M_m$ in $U_3$ is itself an instance of the reduction principle
of Doc.~296: a \emph{projection onto a limited temporal segment}, just as a
scale projection is a limited spatial segment. It leaves out (everything
beyond $M_m$) and thereby makes a partial domain evaluable --- the local,
most recent dynamics. This is the same mechanism as in Doc.~296 (reduction
as the condition of evaluability), only applied to the time axis rather than
the scale axis. The window acts in the same sense as any projection.

The decisive difference therefore lies not in the fact \emph{that} $U_3$
reduces --- every evaluable description reduces --- but in \emph{whether the
whole is named of which the reduction is a segment}.

FFGFT's reductions are exhibited as reductions of a known whole: the frozen
75-turn winding and the Markov limit of the memory kernel are segments of
the full, accumulating, scale-invariant structure --- one knows what has
been left out and how to recover it (let the recursion run, do not freeze
$\xi$, Doc.~295/283). The way back is laid out.

Marcel's window, by contrast, stands on its own, without the full object
behind it --- it is the projection without the projected. Recovery is not
laid out because the full object is not named. The window evaluates the
local temporal structure --- a legitimate focus --- but it is not
recognisable as a segment of a boundaryless whole (Doc.~296: the embedded
observer who sits on their scale without unrolling the recursion).

This is precisely what makes the embedding $\iota$ natural: it supplies the
missing whole. $U_3$ as a reduction receives back, in FFGFT, the object of
which it is a segment. The window acts in the same sense as any projection
--- the only difference is that FFGFT knows the way back and HLV does not
(so far) lay it out.

\subsection*{Pre-registered prediction for the OP8-Neuro test (6 July 2026)}

Marcel Krueger has published, with HLV-CONSCIOUSNESS-BRIDGE1 v0.3a
(OP8-LOCK), a data-freeze / protocol-lock that fixes --- before the run ---
a model ladder and a null family:
\begin{center}
\begin{tabular}{ll}
$M_0$ & $\Delta\Phi$ / S-I-C only \\
$M_1$ & $U_1$ only \\
$M_2$ & $U_1 + U_2$ \\
$M_3$ & $U_1 + U_2 + U_3$ \\
$M_4$ & generic ML baseline
\end{tabular}
\end{center}
Central question (his formulation): whether the full $U_1/U_2/U_3$ model
($M_3$) provides measurable additional value over $U_1/U_2$, information-only,
coherence-only, \textbf{memory-free Markov}, phase-shuffled, time-shuffled,
and the generic ML baseline --- the \emph{$U_3$ necessity test}.

This pre-registered test is simultaneously the decisive test for the
embedding hypothesis of this document. FFGFT makes the following prediction
about it --- \emph{before} Marcel's run, explicitly booked as a prediction
(P35), not as a claim on his result:

\begin{enumerate}
  \item \textbf{If $U_3$ is a hard-reset window with $M_m < 75$ (winding
    coordinate)} --- i.e.\ a segment below one full turn, the Markov limit
    of the embedding $\iota$ --- then \textbf{$M_3$ does not significantly
    outperform the memory-free Markov baseline}. $U_3$ then lies entirely in
    the degenerate boundary case; the additional channel carries no
    structure beyond Markov.
  \item \textbf{If $M_3$ clearly outperforms the memory-free Markov
    baseline}, then $U_3$ carries non-Markov structure (correlations beyond
    $M_m$, effective accumulation) that the hard-reset kernel does not
    explain. In that case the simple embedding $\iota$ into the Markov limit
    is \emph{falsified}, and a richer embedding must be constructed.
\end{enumerate}

The prediction is sharp because it hangs on Marcel's own construction: a
bounded window $M_m$ of fixed width (Eq.~7 of his Consciousness document,
\enquote{avoids uncontrolled accumulation}) can by construction carry no
correlations beyond $M_m$ --- so the expected outcome is (1). An outcome (2)
would indicate that $U_3$, in implementation, does more than its definition
suggests --- which would be independently interesting and would then need to
be localised.

This prediction does not touch Marcel's claim boundaries: it says nothing
about consciousness, no physical validation, no HLV ontology. It concerns
solely the structural question of whether $U_3$ carries beyond the Markov
limit --- and that is exactly what his OP8 test probes.

\subsection*{Status}

Candidate note. Directed speculation on the basis of secured partial
findings (P35). The subset conjecture is structurally plausible and
falsifiable --- that is its value. It is not a claim.
